% ---------------------------------------------------------------- Background
\section*{Background}
Training and fine-tuning jobs keep consuming compute long after they stop improving~\cite{schwartz2020greenai}. In practice they are supervised either by a person reading loss curves or by an automatic rule: patience on a validation signal~\cite{prechelt2012early}, the median stopping rule~\cite{golovin2017vizier}, or multi-fidelity schedulers such as successive halving, Hyperband~\cite{li2018hyperband}, ASHA~\cite{li2020asha} and BOHB~\cite{falkner2018bohb}, available in Optuna~\cite{akiba2019optuna}, Ray Tune~\cite{liaw2018tune} and the Kubernetes-native Katib~\cite{george2020katib}. Learning-curve extrapolation predicts a run's final performance from a partial curve with parametric ensembles~\cite{domhan2015extrapolation}, Bayesian neural networks~\cite{klein2017lcnet} or prior-data fitted networks~\cite{adriaensen2023lcpfn}. These methods save compute, but the reason a particular job was stopped is a rank, a quantile or a posterior rather than a statement an operator can check. Large language model (LLM) agents are now applied to the same task~\cite{zhang2023llmhpo,liu2024agenthpo,huang2024mlagentbench,chan2025mlebench}; they add natural-language rationales, but nothing guarantees that a rationale is what produced the action~\cite{yao2023react}.

For consequential automated decisions, models that are interpretable by construction are preferable to black boxes explained after the fact~\cite{rudin2019stop}. Kolmogorov--Arnold Networks (KANs)~\cite{liu2025kan} place learnable univariate functions on edges and can be converted into symbolic formulas~\cite{liu2024kan2}, which makes them a candidate for a controller whose decision rule can be read. Symbolic models have been used to explain hyperparameter optimisation after the fact~\cite{segel2023symbolic} and to distil learned optimisers~\cite{zheng2022symbolicl2o}, and KANs have produced interpretable reinforcement-learning policies~\cite{singh2025kanrl}. A controlled comparison finds that multilayer perceptrons generally match or exceed KANs except on symbolic-formula tasks~\cite{yu2024kanormlp}; we therefore claim no accuracy advantage for the KAN and compare it with linear, tree and MLP gates trained on identical labels. Placing a learned symbolic rule between an agent and the actions it may take also resembles shielding in safe reinforcement learning~\cite{alshiekh2018shielding}.

The method described here customises the original KAN method in two ways. First, it supplies a supervision signal for a control decision that has no ground truth at decision time: hindsight labels derived from completed runs (Step 3). Second, it turns the trained KAN into an executable closed-form decision rule and checks that the rule and the network agree (Step 4). It also adds a deployment guard against a failure mode we observed in closed-loop control (Step 5). The method was developed inside an agentic controller for Kubernetes training jobs, \texttt{KAN-KubeAgent}, but it does not depend on Kubernetes, on an LLM or on any particular training framework: it needs only per-epoch validation loss and gradient norm.

\paragraph{What was validated, and how}
We collected \nRuns{} complete, uncontrolled training runs (\nEpochsRecorded{} epochs) on four real workloads that differ in dataset, architecture and optimiser, and evaluated the gate against eleven alternatives with two protocols: cross-validation within each workload, grouped by seed, and leave-one-workload-out transfer, in which the gate controls a workload it has never seen. Uncertainty is reported with run-level bootstrap confidence intervals~\cite{efron1993bootstrap} and differences are tested with paired Wilcoxon signed-rank tests~\cite{wilcoxon1945} with Holm correction~\cite{holm1979}. We then let the gate control live training processes, including real learning-rate changes, on all four workloads. No synthetic or simulated data enters any result.
